\section{Introducción}

El presente informe tiene por objetivo el diseño sísmico de los elementos estructurales del edificio
en estudio: muros de hormigón armado Pier~C y Pier~J, y vigas sísmicas B7, B8 y B39. El análisis
se enmarca en el proyecto de un edificio de hormigón armado de 21 pisos más un nivel de subterráneo,
ubicado en Antofagasta, zona sísmica~3, suelo tipo~A, con categoría de ocupación~II.

En primer lugar, se diseñan los muros sísmicos Pier~C ($l_w = 5{,}40\;\text{m}$) y Pier~J
($l_w = 5{,}76\;\text{m}$), ambos con espesor $e_w = 25\;\text{cm}$ y altura total
$H_w = 50{,}82\;\text{m}$. Para cada piso representativo se determinan las armaduras de corte
horizontal con sobrerresistencia $\Omega_v = 2{,}0$, y las armaduras de flexo-compresión mediante
el método de Cárdenas-Magura. Se verifican además las condiciones de compresión máxima, esbeltez
mínima y corte máximo admisible. A partir del desplazamiento de demanda sísmica calculado según
NCh433 y DS~60, se obtienen la curvatura última y la longitud de compresión requerida $C_c$, que
determinan la necesidad de confinamiento especial de borde en cada nivel.

En segundo lugar, se diseñan las vigas sísmicas B7, B8 y B39 a lo largo de los 20 pisos del
edificio, determinando su armadura longitudinal a flexión y transversal a corte a partir de la
envolvente de esfuerzos obtenida mediante modelación en SAP2000, conforme a ACI~318-19.